\section{Learning Channel Denoising and Activity Detection}
\label{sec:learning}

This section develops learning methods for the conditional channel moments
and activity probability required by Section~\ref{sec:joint-denoising}.
The channel moments can be expressed through score functions of the
noisy-channel distribution.
We use denoising score matching to learn a channel denoiser and obtain
the variance in \eqref{eq:conditional-channel-variance} from its input
derivatives \cite{dytso2023,manor2024}.
For activity detection, a classifier estimates the density ratio in
\eqref{eq:local-activity}.

\subsection{Channel Mean Estimation}
\label{sec:score-denoising}
\label{sec:channel-learning}

Under the active hypothesis in \eqref{eq:pseudo-awgn}, the complex-valued
channel $\mathbf H\sim q_{\rm tr}$ is corrupted by proper
complex Gaussian noise with covariance $v\I_D$, where $v>0$.
The noise is independent of $\mathbf H$.
Stacking the real and imaginary parts of $\vect(\mathbf H)$ as
$\mathbf h_{\mathbb R}\in\mathbb R^{2D}$ gives the real-valued
observation model
\begin{equation}
 \mathbf u=\mathbf h_{\mathbb R}+\sqrt{v/2}\,\bm\epsilon,
 \qquad \bm\epsilon\sim\mathcal N(\mathbf0,\I_{2D}).
 \label{eq:training-query}
\end{equation}

At fixed $v$, let $q_v^{\mathbb R}$ denote the density of
$\mathbf u$.
The conditional mean is given by Tweedie's formula
\cite{efron2011,dytso2023}
\begin{equation}
 \mathbf m_{\mathbb R}(\mathbf u,v)
 =\mathbf u+\frac v2\nabla_{\mathbf u}
       \log q_v^{\mathbb R}(\mathbf u),
 \label{eq:tweedie-mean}
\end{equation}
where $\nabla_{\mathbf u}\log q_v^{\mathbb R}$ is the
first-order score function \cite{hyvarinen2005}.

Let $\mathbf s_\theta(\mathbf u,v)$ approximate the first-order score, with
trainable parameters $\theta$.
Although the score of $q_v^{\mathbb R}$ is unknown, the conditional
Gaussian density in \eqref{eq:training-query} has score
$-2(\mathbf u-\mathbf h_{\mathbb R})/v$.
Denoising score matching \cite{vincent2011} therefore uses
\begin{equation}
 \mathcal L_{\rm DSM}(\theta)=
 \E\left[\frac{v}{4D}
 \norm{\mathbf s_\theta(\mathbf u,v)
 +\frac2v(\mathbf u-\mathbf h_{\mathbb R})}^2\right].
 \label{eq:denoising-score-objective}
\end{equation}
For each fixed $v$, this loss has the same minimizers as the
mean-squared error between the network output and the score of the
perturbed channel distribution.

To obtain the channel estimate directly, we parameterize the score as
$\mathbf s_\theta=2(\mathbf f_\theta-\mathbf u)/v$, where $\mathbf f_\theta(\mathbf
u,v)$ is a denoiser.
Substituting this relation into \eqref{eq:denoising-score-objective}
gives the training loss
\begin{equation}
 \mathcal L_{\rm ch}(\theta)=
 \E\left[\frac{1}{Dv}
 \norm{\mathbf f_\theta(\mathbf u,v)-\mathbf h_{\mathbb R}}^2\right].
 \label{eq:first-loss}
\end{equation}
The expectation is taken over the clean channels, noise variances, and
Gaussian noise.
For each fixed $v$, the conditional mean $\mathbf m_{\mathbb R}(\mathbf
u,v)$ minimizes this loss over all estimators of $\mathbf h_{\mathbb R}$ from $\mathbf u$.

The denoiser is trained over a range of noise variances
\cite[Sec.~IV-B]{cai2026} to accommodate the input messages in
\eqref{eq:pseudo-awgn}. Its real and imaginary outputs are reshaped into
the complex-valued estimate $\mathbf F_\theta(\mathbf R,v)$.
The sampling distribution of $v$ is specified in
Section~\ref{sec:training-settings}, and the architecture and
preconditioning of $\mathbf f_\theta$ are given in
Section~\ref{sec:network-settings}.

\subsection{Channel Variance Estimation}
\label{sec:conditional-error}

The conditional covariance $\mathbf P_{\mathbb
R}=\operatorname{Cov}(\mathbf h_{\mathbb R} \mid\mathbf u,v)$ is
related to the second-order score \cite{dytso2023},
\cite[Sec.~IV-A]{cai2026}:
\begin{equation}
 \mathbf P_{\mathbb R}
 =\frac v2\I_{2D}+\frac{v^2}{4}
   \nabla_{\mathbf u}^2\log q_v^{\mathbb R}(\mathbf u).
 \label{eq:tweedie-covariance}
\end{equation}
Here the second-order score is the Hessian of the log density.

Differentiating \eqref{eq:tweedie-mean} with respect to $\mathbf u$
at fixed $v$ and comparing the result with \eqref{eq:tweedie-covariance}
relates the covariance to the mean Jacobian \cite[Prop.~1]{dytso2023}.
Since the joint denoiser uses the average variance in
\eqref{eq:conditional-channel-variance}, only the covariance trace is
needed:
\begin{align}
 \mathbf P_{\mathbb R}&=\frac v2\mathbf J_m,
 \label{eq:mean-derivative-covariance}\\
 \tau^\star&=\frac{v}{2D}\operatorname{div}\mathbf m_{\mathbb R}.
 \label{eq:mean-derivative-variance}
\end{align}
Here $\mathbf J_m=\partial\mathbf m_{\mathbb R}/\partial\mathbf u$ is
the Jacobian at fixed $v$, and $\tau^\star$ is the average
conditional variance per complex coordinate.
Thus, the average variance can be estimated from the derivatives of the
trained denoiser without training a separate variance model
\cite[Sec.~3.2 and Sec.~4.1]{manor2024}.

For numerical evaluation, we approximate the divergence by Monte Carlo
derivative estimation \cite{metzler2016}, using central finite differences
of $\mathbf f_\theta$ to estimate $\tau^\star$.
Let $\mathbf z_\ell$ be independent random sign vectors with $\E[\mathbf
z_\ell\mathbf z_\ell^{\mathsf T}]=\I_{2D}$, and define $\mathbf
w_\ell=\sqrt{v/2}\,\mathbf z_\ell$.
Using $L_z$ random directions and a finite-difference step
$\delta_{\rm FD}>0$, the resulting estimate is
\begin{equation}
 \widehat\tau_\theta(\mathbf R,v)
 =\frac{1}{2\delta_{\rm FD} D L_z}\sum_{\ell=1}^{L_z}
       \mathbf w_\ell^{\mathsf T}\Delta_\ell \mathbf f_\theta,
 \label{eq:derivative-trace}
\end{equation}
where $\Delta_\ell \mathbf f_\theta
 =\mathbf f_\theta(\mathbf u+\delta_{\rm FD}\mathbf w_\ell,v)
 -\mathbf f_\theta(\mathbf u-\delta_{\rm FD}\mathbf w_\ell,v)$.
The finite-difference settings are given in Section~\ref{sec:receiver-settings}.

\subsection{Activity Probability Estimation}
\label{sec:density-ratio}

The activity probability depends on the likelihoods through their ratio,
as seen by writing \eqref{eq:local-activity} in log-odds form:
\begin{equation}
 \log\frac{p^\star(\mathbf R,v)}{1-p^\star(\mathbf R,v)}
 =\log\frac{\rho}{1-\rho}
  +\log\frac{q_{1,v}(\mathbf R)}{q_{0,v}(\mathbf R)}.
 \label{eq:activity-log-odds}
\end{equation}
Direct evaluation of $q_{1,v}$ requires the channel integral in
\eqref{eq:query-densities}. Binary classification instead allows the
ratio $q_{1,v}/q_{0,v}$ to be estimated from observations drawn under
the two hypotheses \cite{menon2016}.

The two hypotheses define the training classes $a\in\{0,1\}$, with
equal prior probability $\Pr(a=1)=1/2$.
The active class $a=1$ consists of observations
$\mathbf R=\mathbf H+\mathbf W$ with $\mathbf H\sim q_{\rm tr}$,
whereas the inactive class $a=0$ contains noise-only observations
$\mathbf R=\mathbf W$.
The noise satisfies $\vect(\mathbf W)\sim\CN(\mathbf0,v\I_D)$ independently of $\mathbf H$,
and both classes use the same sampling distribution of $v$.
We train the classifier with parameters $\phi$ by minimizing the
binary cross-entropy loss
\begin{equation}
 \mathcal L_{\rm act}(\phi)=
 \E[\log(1+e^{z_\phi})-a z_\phi],
 \label{eq:classifier-loss}
\end{equation}
where $z_\phi=z_\phi(\mathbf R,v,\mathbf F_\theta(\mathbf R,v))$.
For this balanced training distribution, the minimizer of the expected
loss over all functions of $(\mathbf R,v)$ is the log likelihood ratio
\begin{equation}
 z^\star(\mathbf R,v)=
       \log\frac{q_{1,v}(\mathbf R)}{q_{0,v}(\mathbf R)}.
 \label{eq:ideal-logit}
\end{equation}

The classifier output therefore estimates the log density ratio
\cite[Sec.~4.2, eq.~(17)]{menon2016}. For
$\mathbf F_k=\mathbf F_\theta(\mathbf R_k,v)$ and
$\ell_k=z_\phi(\mathbf R_k,v,\mathbf F_k)$, combining this estimate
with the receiver's activity prior in \eqref{eq:activity-log-odds} gives
\begin{equation}
 p_k=\sigmoid(\logit\rho+\ell_k),
 \label{eq:activity-probability}
\end{equation}
where $\logit a=\log[a/(1-a)]$ and $\sigmoid b=(1+e^{-b})^{-1}$.

